\scanpage{36}
\mspage{1}
$\varepsilon_0,\varepsilon_1\in H$ (\inserted{\unsure{gén. du}} groupe $H$) tels que
\[\boxed{\varepsilon_0\varepsilon_1\varepsilon_0=\varepsilon_1\varepsilon_0\varepsilon_1}\qquad(*)\]
\margin{\deleted{$\unsure{\sigma^{-1}\rho_0}$}$\unsure{\varepsilon_0^{-1}\varepsilon_1^{-1}\varepsilon_0^{-1}\varepsilon_1}$}
On pose
\[
\left\{\begin{aligned}
\sigma&=\varepsilon_0\varepsilon_1\varepsilon_0\ (=\varepsilon_1\varepsilon_0\varepsilon_1)\\
\rho'_0&=\sigma\varepsilon_0=\deleted{\unsure{\varepsilon_1\sigma}}\ (= (\varepsilon_1\varepsilon_0)^{\unsure{2}})\\
\rho'_1&=\sigma\varepsilon_1=\varepsilon_0\sigma\quad\deleted{\unsure{\varepsilon_0\varepsilon_1\varepsilon_0\varepsilon_1}}\ (=(\varepsilon_0\varepsilon_1)^2)\\
\rho_0&=\sigma^{-1}\varepsilon_0=\varepsilon_1\sigma^{-1}\ (=\varepsilon_0^{-1}\varepsilon_1^{-1}\deleted{\unsure{\varepsilon_0}}=(\varepsilon_1\varepsilon_0)^{-1})\\
\rho_1&=\varepsilon_0\sigma^{-1}=\sigma^{-1}\varepsilon_1\ (=\varepsilon_1^{-1}\varepsilon_0^{-1}=(\varepsilon_0\varepsilon_1)^{-1})
\end{aligned}\right.
\]
\margin{[donc $\sigma(\varepsilon_0)=\varepsilon_1,\sigma(\varepsilon_1)=\varepsilon_0,\ \sigma^2\in\unsure{\operatorname{Centr}}(H)$\par
$\operatorname{int}(\sigma)=\operatorname{int}(\sigma^{-1})$]}
\diagram{$\begin{array}{ccc}\rho_0,\rho'_0&&\rho_1,\rho'_1\\&\sigma&\\\varepsilon_0&&\varepsilon_1\end{array}$}
Alors\quad\deleted{$\sigma=\varepsilon_0\rho_0$}
\[
\text{(a)}\quad\left\{\begin{gathered}
\boxed{\begin{gathered}\sigma(\varepsilon_0)=\varepsilon_1,\quad\sigma(\varepsilon_1)=\varepsilon_0\\
\sigma(\rho_0)=\rho_1,\quad\sigma(\rho_1)=\rho_0\end{gathered}}\\
\sigma(\rho'_0)=\rho'_1,\quad\sigma(\rho'_1)=\unsure{\rho'_1|\rho'_0}
\end{gathered}\right.
\]
\margin{conj. par $\sigma$}
\margin{(\uline{NB} \deleted{\unsure{chacune}}\inserted{La première} de ces relations équivaut à $(*)$, et elle implique \inserted{la deuxième}\deleted{\unsure{que}} $\sigma^2$ commute à $\varepsilon_0,\varepsilon_1$, \unsure{donc au groupe engendré} donc (si $H$ engendré par $\varepsilon_0,\varepsilon_1$) $\forall h\in H$, on a aussi $\sigma(h)=\sigma^{-1}(h)$}
\[
\text{(a')}\quad\left\{\begin{gathered}
\rho_0(\varepsilon_0)=\rho'_0(\varepsilon_0)=\varepsilon_1\\
\rho_1(\varepsilon_1)=\rho'_1(\varepsilon_1)=\varepsilon_0\\[2pt]
\varepsilon_0(\rho_0)=\rho_1,\quad\varepsilon_0(\rho'_0)=\rho'_1\\
\varepsilon_1(\rho_1)=\rho_0,\quad\varepsilon_1(\rho'_1)=\rho'_0
\end{gathered}\right.
\]
\margin{conj. par $\rho_i,\rho'_i$\par et par $\varepsilon_i$}
\[
\text{(b)}\quad\begin{gathered}
\rho_0=\sigma^{-1}\varepsilon_0=\varepsilon_1\sigma^{-1},\quad\rho_1=\sigma^{-1}\varepsilon_1=\varepsilon_0\sigma^{-1}\\
\rho'_0=\sigma\varepsilon_0=\varepsilon_1\sigma,\quad\rho'_1=\sigma\varepsilon_1=\varepsilon_0\sigma
\end{gathered}
\]
\margin{$\rho_i,\rho'_i$ en termes de $\sigma,\varepsilon_i$}
\deleted{$\text{(c)}\quad\unsure{\rho_0^3=\rho_1^3=\sigma^{-2}}\quad\rho'_0=\rho_0^{-2},\quad\rho'_1=\rho_1^{-2}$}
\margin{\deleted{(car $\rho_0'^{-1}\rho_0=\rho_0^3$ commute à $\varepsilon_0$, \unsure{donc} $\rho_1'^{-1}\rho_1=\rho_1^3$ commute à $\varepsilon_1$)}}
\[
\text{(c)}\quad\rho_0^{-3}=\rho_1^{-3}=\sigma^2\quad(\mathrel{\mathop=^{\rm def}}\deleted{\unsure{\varepsilon}}\omega\in Z(H))
\]
\[
\text{(d)}\quad\rho'_0=\deleted{\unsure{\sigma}}\rho_0^{-2}=\omega\rho_0,\quad\rho'_1=\rho_1^{-2}=\omega\rho_1
\]
\deleted{$\unsure{\rho'_0}$}
\margin{(donc $\operatorname{int}(\rho_0)=\operatorname{int}(\rho'_0)=\operatorname{int}(\rho_0)^{-2}$)}
\margin{d'où $\rho_0=\omega^{-1}\rho'_0$\par$(=\rho'_0\omega^{-1})$,\quad$\rho_1=\omega^{-1}\rho'_1=\rho'_1\omega^{-1}$}

\uline{NB} Il vaut mieux au début ne pas introduire $\rho'_0,\rho'_1$, quitte à \unsure{la fin de} donner un formulaire complémentaire les concernant.
\margin{(\unsure{Juin} 1983)}
\mspage{35}
